\documentclass[11pt]{article}
\usepackage[margin=1in]{geometry}
\usepackage{amsmath,amssymb,amsthm}
\usepackage[colorlinks=true,urlcolor=blue,linkcolor=blue,citecolor=blue]{hyperref}
\newtheorem{theorem}{Theorem}[section]
\newtheorem{lemma}[theorem]{Lemma}
\newtheorem{corollary}[theorem]{Corollary}
\newtheorem{remark}[theorem]{Remark}
\newcommand{\tr}{\operatorname{tr}}
\newcommand{\rank}{\operatorname{rank}}
\newcommand{\diag}{\operatorname{diag}}
\newcommand{\N}{\mathcal N}
\newcommand{\R}{\mathbb R}
\title{Relative fractional edge updates on unweighted graphs:\\
the degree-two versus degree-three dichotomy}
\author{Independent unweighted research extension}
\date{November 2026}
\begin{document}
\maketitle
\begin{abstract}
Fix $0<\alpha<1$. For a simple unweighted host graph, double one existing
edge and measure the fractional-Laplacian correction relative to the
original fractional energy. Uniformly over all unit-weight paths and
edge locations, its optimal compression rank at tolerance $\varepsilon$
is $\Theta_\alpha(\log^2(1/\varepsilon))$. The same worst-case order holds
for all connected graphs of maximum degree two. The centered infinite-path
limit has the precise count constant $1/(2\pi^2)$, independently of
$\alpha$; this follows by reduction to Laptev's classical weighted
Cauchy-kernel theorem. In contrast, simple unweighted graphs of maximum
degree three already have worst-case rank
$\Theta_\alpha(1/\varepsilon)$. A binary fanout, long terminal necks,
and separated expander-reservoir masses give a finite existence
construction, which can moreover be chosen bipartite.
We prove the reservoir limit in the relative metric,
including the surviving contact resistances. At every fixed maximum
degree there is also a compulsory positive relative outlier, precluding
the full determinant-admissible spectral closure available for arbitrary
weighted paths. Sharp uniform leading constants and historical priority
of these graph-relative specializations are not claimed.
\end{abstract}

\section{Question and results}
All graphs are finite, connected and simple, with unit vertex masses
and unit edge conductances before the update. If $b=e_u-e_v$ is an
existing edge, write $L'=L+bb^T$ and
\[
 C_\alpha=(L^\dagger)^{\alpha/2}(L'^\alpha-L^\alpha)
                                      (L^\dagger)^{\alpha/2},\qquad
 \N_C(\varepsilon)=\#\{j:\lambda_j(C)>\varepsilon\}.
\]
Every inverse power is zero on the common constant kernel. Let
$W_{\mathrm{path},\alpha}$ be the supremal strict count over paths,
and $W_{d,\alpha}$ the supremal count over hosts of maximum degree at
most $d$. An edge conductance changing from one to two does not change
the combinatorial degree.

\begin{theorem}[Worst-case orders]\label{thm:main}
As $\varepsilon\downarrow0$, for each fixed $0<\alpha<1$,
\[
 W_{\mathrm{path},\alpha}(\varepsilon)
   =\Theta_\alpha(\log^2(1/\varepsilon)),\qquad
 W_{2,\alpha}(\varepsilon)
   =\Theta_\alpha(\log^2(1/\varepsilon)),
\]
whereas, for every fixed integer $d\ge3$,
\[
 W_{d,\alpha}(\varepsilon)=\Theta_\alpha(1/\varepsilon).
\]
The degree-three lower examples edit a bridge, hence have leverage
exactly one. Explicitly, for $L=\log(1/\varepsilon)$,
\[
 \frac1{2\pi^2}
 \le\liminf_{L\to\infty}\frac{W_{\mathrm{path},\alpha}(e^{-L})}{L^2}
 \le\limsup_{L\to\infty}\frac{W_{2,\alpha}(e^{-L})}{L^2}
 \le\frac{(1+2\alpha)^2}{4\pi^2\alpha}.
\]
The last coefficient is a nonsharp upper bound.
\end{theorem}

The strict count is precisely the minimum rank needed to approximate
the positive relative matrix in operator norm to tolerance
$\varepsilon$. A positive spectral truncation corresponds to
$0\preceq (L'^\alpha-L^\alpha)-F\preceq\varepsilon L^\alpha$
on the common support. These are structural existence results;
no running-time or finite-precision bound is asserted.

\section{Resolvents, leverage measure, and the path upper bound}
For $A\succ0$, $B=A+uu^T$, set $\tau=u^TA^{-1}u$ and
$c_\alpha=\sin(\pi\alpha)/\pi$. Sherman--Morrison and the scalar
fractional resolvent integral give
\begin{equation}\label{eq:res}
 C=c_\alpha\int_0^\infty
 \frac{t^\alpha v_tv_t^T}{1+m(t)}\,dt,\quad
 v_t=A^{-\alpha/2}(A+tI)^{-1}u,\quad
 m(t)=u^T(A+tI)^{-1}u.
\end{equation}
Write $H$ for the same integral without its denominator. Since
$0\le m(t)\le\tau$,
\begin{equation}\label{eq:sandwich}
 (1+\tau)^{-1}H\preceq C\preceq H.
\end{equation}
In an eigenbasis $A\phi_k=\lambda_k\phi_k$, put
$z_k=|\langle u,\phi_k\rangle|^2/\lambda_k$ and
$\nu=\sum_kz_k\delta_{\lambda_k}$. Then
\begin{equation}\label{eq:diag}
 H_{kk}=\alpha z_k,\qquad
 \tr(1_{(0,\delta]}(A)C1_{(0,\delta]}(A))
                                      \le\alpha\nu((0,\delta]).
\end{equation}
The diagonal identity follows by differentiating the scalar power
integral. Also, after diagonal sign conjugation,
\begin{equation}\label{eq:kernel}
 H_{ij}=\sqrt{z_iz_j}K_\alpha(\log\lambda_i-\log\lambda_j),\qquad
 K_\alpha(s)=\frac{\sinh(\alpha s/2)}{\sinh(s/2)},\quad K_\alpha(0)=\alpha.
\end{equation}

\begin{lemma}[Exact unit-path measure]\label{lem:pathmeasure}
On the path with $n$ vertices, for the edge $(p,p+1)$,
\[
 \lambda_k=4\sin^2\frac{k\pi}{2n},\qquad
 z_k=\frac2n\sin^2\frac{k\pi p}{n},\quad 1\le k<n.
\]
In particular $\nu((0,\delta])\le\sqrt\delta$, uniformly in $n,p$.
On the cycle with $n\ge3$ vertices, the complex Fourier eigenbasis gives
$\lambda_k=4\sin^2(\pi k/n)$ and $z_k=1/n$, $1\le k<n$;
there $\nu((0,\delta])\le\sqrt\delta/2$.
\end{lemma}
\begin{proof}
The normalized nonconstant path eigenvectors are
$\phi_k(j)=\sqrt{2/n}\cos(k\pi(j-1/2)/n)$.
Subtract their values across the edge and divide the squared difference
by $\lambda_k$. Since $\sin x\ge2x/\pi$ for $0\le x\le\pi/2$,
$\lambda_k\ge4k^2/n^2$. Thus at most $n\sqrt\delta/2$
positive eigenvalues are at most $\delta$, each having leverage weight
at most $2/n$. On a cycle,
$|b^T\phi_k|^2=\lambda_k/n$ and
$\lambda_k\ge16\min(k,n-k)^2/n^2$; counting both ends of the
Fourier grid proves the stated bound. Complex eigenvectors merely
diagonalize a real self-adjoint operator and do not change its spectrum.
\end{proof}

\begin{lemma}[Positive pinching]\label{lem:pinch}
For complementary orthogonal projections $P,Q$ and $C\succeq0$,
$C\preceq2(PCP+QCQ)$. In particular, if
$\|QCQ\|\le\varepsilon/2$, then
$\N_C(\varepsilon)\le\N_{PCP}(\varepsilon/2)$.
\end{lemma}
\begin{proof}
Apply $2\langle C^{1/2}Px,C^{1/2}Qx\rangle
\le\|C^{1/2}Px\|^2+\|C^{1/2}Qx\|^2$.
The upper bound is block diagonal; its $Q$ block has no eigenvalue
strictly greater than $\varepsilon$.
\end{proof}

\begin{proof}[Uniform upper bound in Theorem~\ref{thm:main}]
Both endpoint spectra lie in $[0,6]$ for paths and cycles. Take
$\delta=(\varepsilon/(2\alpha))^2$,
$Q=1_{(0,\delta]}(L)$, and $P=I-Q$ on the nonconstant support.
Equations~\eqref{eq:diag} and Lemma~\ref{lem:pathmeasure} give
$\|QCQ\|\le\varepsilon/2$.
Stahl's Theorem~1~\cite{Stahl} supplies a real rational function $r_q$
of type $(q,q)$, without poles on $[0,6]$, such that
\[
 \sup_{[0,6]}|r_q(x)-x^\alpha|
                         \le K_\alpha e^{-2\pi\sqrt{\alpha q}}.
\]
The classical rational rank identity is
$\rank(r_q(L')-r_q(L))\le q$.
For completeness, write $r_q=a_q/d_q$ after canceling common factors,
and factor
\[
 a_q(x)d_q(y)-d_q(x)a_q(y)
           =(x-y)\sum_{i,j<q}c_{ij}x^iy^j.
\]
Left-right evaluation at $L',L$ shows that the range after multiplying
by the invertible denominators is contained in
$\operatorname{span}\{(L')^ib:0\le i<q\}$, proving the rank statement.
It follows that $PCP$ has a rank-$q$ approximant at error at most
$2K_\alpha\delta^{-\alpha}e^{-2\pi\sqrt{\alpha q}}$.
Choose this error at most $\varepsilon/2$. With
$L_\varepsilon=\log(1/\varepsilon)$ it is enough to take
\[
 q=\frac{[(1+2\alpha)L_\varepsilon+O_\alpha(1)]^2}{4\pi^2\alpha}+O(1).
\]
Lemma~\ref{lem:pinch} proves the asserted uniform count bound.
Every connected graph of maximum degree two with an edge is a path or
a cycle, so it is covered as well.
\end{proof}

\begin{remark}
The same proof works for any uniformly bounded-endpoint rank-one family
whose leverage measures satisfy $\nu((0,\delta])\le K\delta^\beta$
with fixed $K,\beta>0$. Its nonsharp coefficient is
$(1+\alpha/\beta)^2/(4\pi^2\alpha)$. Bounded degree alone does not
imply such a small-energy measure bound.
\end{remark}

\section{A matching lower bound from the centered-path limit}
Let $C_N$ be the correction for the central edge of the path on $2N$
vertices. In Lemma~\ref{lem:pathmeasure}, exactly the odd indices
contribute, each with weight $1/N$. Their leverage measures converge to
\begin{equation}\label{eq:arcsine}
 d\nu(\lambda)=\frac{d\lambda}{\pi\sqrt{\lambda(4-\lambda)}},\quad
 0<\lambda<4,
\end{equation}
which is uniform $d\theta/\pi$ under
$\lambda=4\sin^2(\theta/2)$, $0<\theta<\pi$.

\begin{theorem}[Exact centered-limit count]\label{thm:center}
After isometric embeddings of their cyclic supports, $C_N$ converge
in Hilbert--Schmidt norm to a positive compact operator $C_\infty$,
and
\[
 \N_{C_\infty}(\varepsilon)
             \sim\frac1{2\pi^2}\log^2(1/\varepsilon).
\]
The constant is independent of fixed $0<\alpha<1$.
\end{theorem}
\begin{proof}
The limiting derivative is the integral operator on $L^2(\nu)$ with
kernel $K_\alpha(\log\lambda-\log\eta)$.
Identify the finite cyclic space with normalized step functions on the
$N$ equal cells in $(0,\pi)$; the contributing angles are the cell
midpoints. The embedded derivative kernels converge pointwise almost
everywhere and are bounded by $\alpha$. Thus they converge in
Hilbert--Schmidt norm.

The exact kernels on $L^2(\nu_N)$ are
\[
 F_N(\lambda,\eta)=c_\alpha(\lambda\eta)^{(1-\alpha)/2}
 \int_0^\infty\frac{t^\alpha\,dt}
 {(\lambda+t)(\eta+t)(1+m_N(t))},\quad
 m_N(t)=\int\frac{\lambda}{\lambda+t}\,d\nu_N(\lambda).
\]
For $t>0$, midpoint quadrature gives
$m_N(t)\to1-\sqrt{t/(t+4)}$.
At every positive $\lambda,\eta$, dominated convergence in $t$
gives convergence of $F_N$, including when its arguments converge to
positive limits. Moreover $0\le F_N\le K_\alpha\le\alpha$.
Dominated convergence on the angle square proves the asserted
Hilbert--Schmidt limit without using an unbounded inverse power on an
infinite graph. Since every edited edge is a bridge, the limiting
operators satisfy $H_\infty/2\preceq C_\infty\preceq H_\infty$.

With convention $\widehat K(\xi)=\int e^{-i\xi s}K(s)\,ds$, the beta
integral, equivalently NIST DLMF Table~1.14.1~\cite{DLMF}, gives
\begin{equation}\label{eq:symbol}
 \widehat K_\alpha(\xi)=
 \frac{2\pi\sin(\pi\alpha)}{\cosh(2\pi\xi)+\cos(\pi\alpha)}.
\end{equation}
For fixed $\alpha$, this symbol is bounded above and below by positive
constant multiples of its value at $\alpha=1/2$.
Hence the convolution operators, their compressions, and their
conjugations by the same real multiplier satisfy the corresponding
Loewner comparisons.

Set $x=\log\lambda$. The unitary passage to Lebesgue measure turns
$H_\infty$ into the weighted convolution kernel
\[
 \sqrt{r(x)r(y)}K_\alpha(x-y),\qquad
 r(x)=\frac{e^{x/2}}{\pi\sqrt{4-e^x}},\quad x<\log4.
\]
On $x<0$, $r(x)=e^{x/2}v(x)^2$ with $v$ bounded and boundedly
invertible. The low-energy compression therefore has eigenvalues
comparable, with fixed multiplicative constants, to those of
\[
 S(x,y)=e^{(x+y)/4}K_{1/2}(x-y),\qquad x,y<0.
\]
Here multiplier comparison is a min--max comparison of
$M_vSM_v$, not an inference from a pointwise kernel inequality.
As $K_{1/2}(s)=1/(2\cosh(s/4))$, the unitary substitution
$u=e^{x/2}$ changes $S$ into the operator on $L^2(0,1)$ with kernel
$2\sqrt{uv}/(u+v)$.
Laptev's theorem~\cite{Laptev}, Theorem~(3), states that the singular
count for the kernel $u^\beta v^\gamma/(u+v)^a$ is
\[
 \N_T(\varepsilon)\sim
 \frac{\log^2(1/\varepsilon)}{2\pi^2\rho},\qquad
 a>0,\quad\beta,\gamma>-1/2,\quad\rho=\beta+\gamma-a+1>0.
\]
Take $a=1$, $\beta=\gamma=1/2$, and $\rho=1$; the kernel is positive,
so the singular and eigenvalue counts agree. Constant threshold factors
preserve the leading logarithmic-square coefficient.

The high-energy compression, $1\le\lambda\le4$, has count
$O_\alpha(\log(1/\varepsilon))$: its kernel
$K_\alpha(\log\lambda-\log\eta)$ is analytic in $\lambda$ in a fixed
complex neighborhood of $[1,4]$, uniformly for real $\eta\in[1,4]$.
Polynomial approximation in this first variable produces rank $q$
and uniform error $O(e^{-cq})$. Finite measure converts this into the
same operator-norm bound, including its integrable density singularity
at four. Compression min--max gives a lower count from the low block;
Lemma~\ref{lem:pinch} gives an upper count by the low and high diagonal
blocks at half the threshold. Thus $H_\infty$ has the stated coefficient,
as does $C_\infty$ by its factor-two Loewner sandwich.
\end{proof}

Norm convergence implies convergence of each ordered eigenvalue.
Consequently, for every positive strict threshold,
\[
 \sup_N\N_{C_N}(\varepsilon)\ge\N_{C_\infty}(\varepsilon).
\]
This proves the lower half of the degree-two assertions.

\begin{corollary}[An edge near a fixed endpoint]
Fix $p\ge1$, double edge $(p,p+1)$, and let the far endpoint tend to
infinity. The limiting relative correction has count
$\log^2(1/\varepsilon)/(6\pi^2)$ to leading order.
\end{corollary}
\begin{proof}
Its angle measure is $(2/\pi)\sin^2(p\theta)d\theta$.
Its logarithmic density is asymptotic to
$(p^2/\pi)e^{3x/2}$ as $x\to-\infty$.
Below a fixed cutoff preceding its first positive zero, it is comparable
to $e^{3x/2}$. The same symbol comparison and substitution give the
kernel $2u^{3/2}v^{3/2}/(u+v)$, with Laptev parameter $\rho=3$.
The remaining finite-energy interval has the same analytic-kernel
$O(\log(1/\varepsilon))$ upper. Step-kernel convergence is as above,
now incorporating the bounded multiplier $\sqrt2\sin(p\theta)$.
\end{proof}

\section{Separated masses and reverse Schur pivots}
The following elementary finite-dimensional limit is useful for the
degree-three construction.

\begin{lemma}[Separated-mass fractional limit]\label{lem:pivots}
Let $K\succ0$, $K'=K+vv^T$ have dimension $m$, and set
\[
 D_\delta=\diag(\delta^{(m-1)/2},\ldots,1),\quad
 A_\delta=D_\delta K D_\delta,\quad B_\delta=D_\delta K'D_\delta.
\]
For $T_i=\{i,\ldots,m\}$ define
\[
 \kappa_i=\frac{\det K_{T_i}}{\det K_{T_{i+1}}},\quad
 \tau_i=v_{T_i}^TK_{T_i}^{-1}v_{T_i},\quad \tau_{m+1}=0.
\]
Empty determinants are one, and $\kappa_i'$ is defined using $K'$.
In an ordered eigenbasis of $A_\delta$, its relative correction tends
in norm to the diagonal entries
\[
 (\kappa_i'/\kappa_i)^\alpha-1
       =\left(\frac{1+\tau_i}{1+\tau_{i+1}}\right)^\alpha-1.
\]
The result also holds for convergent $K_\delta\to K\succ0$,
$v_\delta\to v$, and is consequently uniform on compact parameter
families contained in this class.
\end{lemma}
\begin{proof}
Order $\lambda_i(A_\delta)$ increasingly. The elementary symmetric
polynomial $e_k$ is the sum of principal minors, so
\[
 e_k(A_\delta)\sim\det K_{T_{m-k+1}}\delta^{k(k-1)/2}.
\]
The index set giving this exponent is unique. If $P_k$ is the product
of the largest $k$ eigenvalues, then $P_k\le e_k\le\binom mk P_k$.
Thus
\[
 \lambda_{m-k}/\lambda_{m-k+1}
 =P_{k+1}P_{k-1}/P_k^2
 \le\binom mk^2 e_{k+1}e_{k-1}/e_k^2=O(\delta).
\]
All adjacent scales separate, making every other product in $e_k$
negligible. Taking successive ratios gives
$\lambda_i\sim\kappa_i\delta^{m-i}$.

Put $u_\delta=D_\delta v$ and let $z_i$ be its spectral leverage
weights. Their sum is the fixed $\tau_1$. For $s>0$,
\[
 m_\delta(\delta^{m-i}s)
 =v^T(K+\delta^{m-i}sD_\delta^{-2})^{-1}v
 \longrightarrow v_{T_i}^T(K_{T_i}+se_ie_i^T)^{-1}v_{T_i}.
\]
Terms preceding $i$ are removed by their divergent diagonal; later
diagonal perturbations vanish. A scalar Schur complement writes the
limit as
$\tau_{i+1}+h_i^2/(\kappa_i+s)$, where
$h_i=v_i-K_{i,T_{i+1}}K_{T_{i+1}}^{-1}v_{T_{i+1}}$ and
$h_i^2/\kappa_i=\tau_i-\tau_{i+1}$.
Along any convergent subsequence of the bounded leverage weights,
their spectral expression has limit
$\sum_{j>i}z_j^*+\kappa_i z_i^*/(\kappa_i+s)$.
Equality for all $s>0$ identifies
$z_i\to\tau_i-\tau_{i+1}$.

In~\eqref{eq:res}, substitute $t=\lambda_i s$ for its $i$th diagonal
entry. The denominator tends to
$1+\tau_{i+1}+(\tau_i-\tau_{i+1})/(1+s)$.
The integrable envelope $s^\alpha/(1+s)^2$, together with bounded
$z_i$, justifies dominated convergence. The scalar fractional formula
then gives the asserted diagonal entry. For $i<j$, dropping the positive
denominator bounds the off-diagonal entry by
\[
 \sqrt{z_i z_j}\,x^{(1-\alpha)/2}\frac{1-x^\alpha}{1-x},\qquad
 x=\lambda_i/\lambda_j\longrightarrow0.
\]
Thus all off-diagonal entries vanish. The determinant lemma on trailing
principal submatrices identifies $\kappa_i'/\kappa_i$ with the displayed
leverage ratio. The same arguments apply along convergent parameter
families. Compactness and a subsequence contradiction give uniformity.
\end{proof}

\section{A degree-three unit-weight lower construction}
Only existence of large cubic expanders is taken from outside this
proof. Bordenave's Theorem~1~\cite{Bordenave}, a proof of Friedman's
theorem, implies that for every sufficiently large even $n$ there is a
simple connected cubic graph with Laplacian gap at least some absolute
$\gamma_0>0$. For example one may take
$\gamma_0=(3-2\sqrt2)/2$.

Subdivide one edge and use its new degree-two vertex as a port.
These reservoirs still have a common positive gap $\gamma$.
Indeed, if $E(x)$ is the new graph energy and $\mu$ the mean on the old
vertices, then the old graph energy is at most $2E(x)$, so its old-vertex
variance is at most $2E(x)/\gamma_0$. The new port deviation obeys
$(x_p-\mu)^2\le2E(x)+4E(x)/\gamma_0$.
Variance about the new mean is smaller. Thus
$\gamma=1/(2+6/\gamma_0)$ is valid. Its port resistance to its mean,
\[
 r=e_p^TL_{\mathrm{res}}^\dagger e_p,
\]
lies in $[0,R]$, where $R=1/\gamma$.
Sizes after subdivision are odd and can approximate any prescribed
positive mass ratio.

\subsection{The fixed fanout circuit}
Let $m=2^k\ge2$, choose an integer $K\ge4(R+1)$, and put $\ell=Km$.
Take a full binary tree with $m$ leaves. Join each leaf to a different
terminal reservoir by a neck of $\ell$ unit edges. Join the root $o$
by one unit edge to an anchor reservoir; this is the edge to double.
Every host degree is at most three, and the edited edge is a bridge.

At zero frequency, eliminate mean-zero coordinates in each reservoir.
The terminal neck resistances become $\ell+r_i$, and the root-anchor
resistance becomes $1+r_0$, or $1/2+r_0$ after doubling. This follows
also by minimizing the reservoir energy at fixed port displacement:
the minimum is its squared displacement divided by $r_i$.
Write
\[
 c_0=(1+r_0)^{-1},\qquad c'_0=(1/2+r_0)^{-1}.
\]
With anchor mean grounded, eliminate the finite core to get SPD terminal
stiffness matrices $K(r)$ and $K'(r)$. Their difference is positive
rank one, as follows from the Schur-complement rank-one formula.

Order terminals arbitrarily. Let $g_j$ be the root admittance, excluding
its anchor branch, when terminals $1,\ldots,j$ are grounded and the
rest are open. Then $g_0=0$ and $0\le g_j\le1/K$.
We claim
\begin{equation}\label{eq:gstep}
 g_j-g_{j-1}\ge\frac1{4Km}.
\end{equation}
Hold the root at potential one before installing the $j$th shunt.
Each installed terminal draws at most $1/\ell$. At depth $h$, the edge
on the new terminal's root-to-leaf path carries at most
$(m/2^h)/\ell=1/(K2^h)$. Its open-circuit leaf potential is therefore
$v_j\ge1-1/K\ge3/4$.
With the root grounded, the leaf's Thevenin resistance $\rho_j$ is
at most the direct path resistance $k$. Adding its terminal neck gives
the exact rank-one admittance identity
\[
 g_j-g_{j-1}=\frac{v_j^2}{\ell+r_j+\rho_j}
                 \ge\frac{9}{32Km}\ge\frac1{4Km},
\]
since $r_j\le R$, $k\le m$ and $\ell+r_j+k\le2Km$.

Put $f(g)=(c'_0+g)/(c_0+g)$. For $0\le r_0\le R$ and $0\le g\le1/K$,
\[
 -\frac{d}{dg}\log f(g)
  =\frac{c'_0-c_0}{(c'_0+g)(c_0+g)}\ge\beta>0,
\]
where one may take
\[
 \beta=\frac{1}{2(R+1/2)(R+1)(2+1/K)(1+1/K)}.
\]
For $T_i=\{i,\ldots,m\}$, Schur complementation through the fixed
core and the scalar determinant lemma at the root give
\begin{equation}\label{eq:detcircuit}
 \frac{\det K'_{T_i}}{\det K_{T_i}}
       =\frac{f(g_{i-1})}{f(g_m)}.
\end{equation}
To check the denominator, the core-only Dirichlet matrix pins all
terminals and has determinant ratio $f(g_m)$. The core-plus-$T_i$
matrix pins just the earlier terminals and has ratio $f(g_{i-1})$.
Their quotient is the terminal Schur-complement determinant ratio.
Consequently every reverse pivot ratio satisfies
\begin{equation}\label{eq:pivotbound}
 \log(\kappa_i'/\kappa_i)
 =\log f(g_{i-1})-\log f(g_i)\ge\frac{h}{m},\qquad h=\frac\beta{4K}>0.
\end{equation}

\subsection{Mass choices and the anchor limit}
Give terminal $i$ mass $M_i=\delta^{-(m-i)}$.
Lemma~\ref{lem:pivots} makes its grounded relative correction converge
to entries at least
$(\kappa_i'/\kappa_i)^\alpha-1\ge\alpha h/m$.
The parameters $r\in[0,R]^{m+1}$ form a compact set; the terminal pairs
are continuous and stay SPD. Therefore one sufficiently small fixed
$\delta$ makes all $m$ eigenvalues at least $3\alpha h/(4m)$,
simultaneously for every allowed contact vector.

Now give the anchor mass $M_0$ and let $M_0\to\infty$, keeping these
terminal masses fixed. Its mass-normalized full stiffness converges
to zero direct-sum the grounded terminal stiffness. There is one zero
eigenvalue throughout, and all other $m$ eigenvalues stay away from
zero in this fixed-dimensional limit. Fractional and pseudoinverse
functional calculus therefore converge. The convergence is again
uniform over the compact contact cube. Choose finite $M_0$ large
enough that the resulting ungrounded mass pair has $m$ relative
eigenvalues at least $5\alpha h/(8m)$.

\subsection{Replacing prescribed masses by ordinary reservoirs}
The following argument is needed because zero-frequency elimination
alone does not commute with fractional powers.

\begin{lemma}[Relative reservoir limit]\label{lem:reservoir}
Fix the finite core. Fix positive masses $M_0,\ldots,M_m$.
Choose uniform-gap reservoirs of sizes $n_i(s)$ with $n_i(s)/s\to M_i$.
Along a subsequence on which their contact resistances converge, the
compression of the graph relative correction to its baseline low
spectral subspace converges to the relative correction of the
mass-normalized zero-frequency reservoir Schur pair.
There are $m+1$ low directions, including the common zero direction.
No convergence of the high relative correction is asserted.
\end{lemma}
\begin{proof}
Let $P_s$ be the isometry of normalized reservoir indicators, and use
its orthogonal complement for all mean-zero reservoir coordinates and
all core coordinates. The Laplacian blocks have the form
\[
 L_s=\begin{pmatrix}A_s&B_s\\B_s^T&F_s\end{pmatrix},\qquad
 \|A_s\|=O(s^{-1}),\quad\|B_s\|=O(s^{-1/2}),\quad F_s\succeq\eta I,
\]
with $\eta>0$ independent of reservoir sizes; the same holds for $L_s'$.
To verify the uniform last inequality, mean-zero reservoir norms are
controlled by their internal energy divided by $\gamma$. Core energy
with zero reservoir pins is bounded by twice the actual external energy
plus twice the sum of squared port deviations. This controls the
fixed-dimensional core norm by its positive Dirichlet gap. Combining
these two bounds gives $F_s\succeq\eta I$.
All full endpoint norms are uniformly bounded, by six and eight.

The static Schur matrix is $S_s=A_s-B_sF_s^{-1}B_s^T$.
The exact energy minimization above identifies $sS_s$ with the
mass-normalized core stiffness with contacts $r_i(s)$, up to the
vanishing difference between $n_i(s)/s$ and $M_i$.
Along contact-convergent subsequences it tends to $H$, and similarly
$sS_s'\to H'$. Both have one common zero and rank $m$.
The frequency-dependent Schur correction differs from the static one
by $O(|\lambda|/s)$ for $|\lambda|<\eta/2$.
Standard block elimination, or the graph-subspace equation
$X_s=-F_s^{-1}B_s^T+O(s^{-3/2})$, now gives a baseline low isometry
$U_s=P_s+O(s^{-1/2})$ with
\[
 sU_s^TL_sU_s\longrightarrow H.
\]
There are exactly $m+1$ eigenvalues of order $s^{-1}$, including zero;
all others are bounded below. The updated isometry $U_s'$ has the same
properties, and their coordinate choices can be aligned to $P_s$.

Their mutual high-space leakage is $O(s^{-1/2})$. Thus compressing
the bounded updated high fractional power to the baseline low space
contributes only $O(s^{-1})$. After multiplication by $s^\alpha$ this
vanishes, since $\alpha<1$. The low part gives
\[
 s^\alpha U_s^T L_s'^\alpha U_s\longrightarrow H'^\alpha,
 \qquad s^\alpha U_s^T L_s^\alpha U_s\longrightarrow H^\alpha.
\]
On the baseline positive low space, the inverse factors scale as
$s^{\alpha/2}$ and converge after that scaling to the negative power
of $H$. Set them to zero on the common kernel. Multiplying these
limits proves convergence of the compressed relative correction.
\end{proof}

Contacts always belong to a compact cube. If the last convergence
were to fail uniformly relative to the contact-dependent limit, a
subsequence of contacts would converge, contradicting
Lemma~\ref{lem:reservoir}. Thus the fixed mass choices in the preceding
subsection, followed by sufficiently large finite reservoir sizes,
give $m$ compression eigenvalues at least $\alpha h/(2m)$.
Compression min--max gives the same lower bound for the full graph.
No spurious assertion that the high response vanishes is needed.

\begin{proof}[Degree-three lower bound in Theorem~\ref{thm:main}]
For small $\varepsilon$, choose a power of two with
$\alpha h/(8\varepsilon)\le m<\alpha h/(4\varepsilon)$.
The finite graph above has at least $m$ eigenvalues strictly exceeding
$\varepsilon$.
Conversely, for every existing unit edge, $\tau=b^TL^\dagger b\le1$.
The determinant lemma on the common support gives
\[
 \sum_j\log(1+\lambda_j(C_\alpha))=\alpha\log(1+\tau)\le\alpha\log2.
\]
Thus $\N_C(\varepsilon)<\alpha\log2/\log(1+\varepsilon)$, proving
the matching order upper for all maximum degrees.
\end{proof}

\begin{corollary}[Bipartite degree-three hosts]\label{cor:bipartite}
The degree-three worst-case order remains
$\Theta_\alpha(1/\varepsilon)$ when hosts are additionally required
to be bipartite. The lower examples still double one existing bridge.
\end{corollary}
\begin{proof}
Choose the simple cubic graphs supplied by Bordenave's Theorem~1
with every nontrivial adjacency eigenvalue in
$[-3+\gamma_0,3-\gamma_0]$.
Their bipartite double covers have adjacency matrices
$\left(\begin{smallmatrix}0&A\\A&0\end{smallmatrix}\right)$ and spectra
$\{\mu_i,-\mu_i\}$.
They are simple cubic graphs; the eigenvalues $3,-3$ are simple and
all others lie in the same displayed interval. Thus the covers are
connected and have Laplacian gap at least $\gamma_0$.
For original sizes $2k$ their sizes are $4k$.

Replace one cover edge by a path of length three, adding two vertices.
Its odd length preserves bipartiteness. Either new vertex is a
degree-two port. If $Q$ is the new energy, old-graph energy is at most
$3Q$. With $\mu$ the mean over old vertices, their variance is at most
$3Q/\gamma_0$. Each new vertex is adjacent to an old endpoint, giving
its squared deviation from $\mu$ at most $2Q+6Q/\gamma_0$.
Variance about the new mean is smaller, so the reservoirs have gap
at least
\[
 \gamma_{\mathrm{bip}}=\frac1{4+15/\gamma_0}>0.
\]
Their available sizes $4k+2$ approximate any fixed positive target
mass ratio after common scaling. Their contact resistances therefore
have a uniform bound $R_{\mathrm{bip}}=1/\gamma_{\mathrm{bip}}$.

Use these reservoirs in the preceding fanout construction, with
$K\ge4(R_{\mathrm{bip}}+1)$. The core is a tree, and each reservoir
is joined to it at only one port. Its bipartition can be independently
flipped to agree across that attachment, so the entire host is
bipartite. All degrees remain at most three. The circuit, separated-mass,
anchor, and relative-reservoir arguments use only the uniform gap,
available sizes, and contact bound. They therefore give the same
$\Omega_\alpha(1/\varepsilon)$ count. The determinant upper bound
applies unchanged.
\end{proof}

\section{A compulsory local outlier at fixed degree}
\begin{theorem}\label{thm:outlier}
For every connected simple unweighted host of maximum degree at most
$d$, with one existing edge doubled,
\[
 \lambda_1(C_\alpha)\ge
             \left(1+\frac{2}{d+1}\right)^\alpha-1=:\zeta_{\alpha,d}>0.
\]
In particular full determinant-admissible spectral closure is impossible
at fixed degree, even after padding spectra by zeros.
\end{theorem}
\begin{proof}
More generally, let $A\succ0$, $B=A+uu^T$, $\eta=\|u\|^2>0$,
and $a=u^TAu/\eta$. Scalar Jensen applied to the spectral probability
measure of $u/\sqrt\eta$ gives $m(t)\ge\eta/(a+t)$.
The function $z\mapsto z^2/(1+z)$ is increasing for $z\ge0$.
Formula~\eqref{eq:res}, evaluated at the relative Rayleigh vector
$A^{\alpha/2}u$, therefore has numerator
\[
 u^T(B^\alpha-A^\alpha)u
 \ge c_\alpha\int_0^\infty
          \frac{t^\alpha\eta^2}{(a+t)(a+\eta+t)}\,dt
 =\eta\bigl[(a+\eta)^\alpha-a^\alpha\bigr].
\]
Concavity of the scalar power on the same spectral probability measure
gives the denominator $u^TA^\alpha u\le\eta a^\alpha$. Consequently
\[
 \lambda_1(C_\alpha(A,u))\ge(1+\eta/a)^\alpha-1.
\]
For an existing unit edge $u=e_p-e_q$, one has $\eta=2$ and
$a=(\deg(p)+\deg(q)+2)/2\le d+1$, proving the assertion.
An arbitrarily long equal-entry determinant-admissible list has all
correction eigenvalues smaller than this compulsory bound.
\end{proof}

For $0<\varepsilon<\zeta_{\alpha,d}$ the determinant budget also implies,
with $q=\N_C(\varepsilon)$,
\[
 (q-1)\log(1+\varepsilon)
 \le\alpha\log(1+\tau)-\alpha\log\left(1+\frac2{d+1}\right).
\]
In particular
\[
 \limsup_{\varepsilon\downarrow0}\varepsilon W_{d,\alpha}(\varepsilon)
       \le\alpha\log\frac{2(d+1)}{d+3}<\alpha\log2.
\]
Thus the full weighted-path leading constant is unattainable at fixed
degree, although the $1/\varepsilon$ order survives for $d\ge3$.

\section{Primary comparison, proof status, and limits}
The analytic ingredients are classical. Stahl's Theorem~1 gives the
root-exponential scalar rational approximation used in the upper.
The rational rank mechanism and rational Krylov matrix-function
updates are established; see Beckermann, Cortinovis, Kressner and
Schweitzer~\cite{BCKS}, especially the rank-one discussion in Section~1,
the exactness in Section~3, and the uniform-approximation comparison
in Theorem~4.5. Those are bounds for the unnormalized matrix-function
update; our relative estimate additionally needs the graph's leverage
measure to control arbitrarily small baseline energies.

Laptev's 1974 Theorem~(3), p.~1039, gives precisely the weighted Cauchy
kernel singular-count asymptotic used here. Formula~\eqref{eq:symbol}
is also tabulated in the DLMF. No new weighted-Hankel asymptotic or
scalar approximation theorem is claimed. The expander input is solely
the existence of cubic reservoirs at every sufficiently large even
size, obtained from Bordenave's Theorem~1 for uniform simple regular
graphs. All cited primary theorem statements and their needed
hypotheses were inspected during this session.

The candidate graph-relative contributions are the unit-path
leverage truncation bound, its matching finite-path lower transfer,
the degree-three fanout/reservoir obstruction, and the compulsory
outlier. Internal parallel proof reviews checked the path upper and
the separated-mass algebra; separate reviews checked the reservoir
argument and primary-source specializations. These are not external
or formal verification, and historical priority has not been certified.
The sharp supremal leading constants for degree two and degree three
remain unproved here. The construction's graph size and reservoir
mass ratios may be enormous; no claim of an efficient extremal example
or construction-time bound is made.

\begin{thebibliography}{9}
\bibitem{Stahl} H.~Stahl. Best uniform rational approximation of
$x^\alpha$ on $[0,1]$. Bull. Amer. Math. Soc. (N.S.) 28 (1993), 116--122.
Theorem~1, equations~(2.1)--(2.2).
\url{https://arxiv.org/abs/math/9301217}.
\bibitem{Laptev} A.~A.~Laptev. Spectral asymptotic behavior of a class of
integral operators. Mathematical Notes 16 (1974), 1038--1043.
Theorem~(3), p.~1039. \url{https://doi.org/10.1007/BF01149794}.
Author-hosted full text:
\url{https://www.ma.ic.ac.uk/~alaptev/Papers/MatZam74.pdf}.
\bibitem{DLMF} NIST Digital Library of Mathematical Functions.
Table~1.14.1, Fourier transforms.
\url{https://dlmf.nist.gov/1.14.T1}.
\bibitem{Bordenave} C.~Bordenave. A new proof of Friedman's second
eigenvalue theorem and its extension to random lifts.
Ann. Sci. Ecole Norm. Sup. 53 (2020), 1393--1439. Theorem~1.
\url{https://arxiv.org/abs/1502.04482}.
\bibitem{BCKS} B.~Beckermann, A.~Cortinovis, D.~Kressner,
M.~Schweitzer. Low-rank updates of matrix functions II: Rational Krylov
methods. SIAM J. Numer. Anal. 59 (2021), 1325--1347.
\url{https://arxiv.org/abs/2008.11501}.
\end{thebibliography}
\end{document}
